\section*{Capítulo \ref{ch:g10:sport-and-health} --- Atividade física e saúde}

\begin{solution}{exo:g10:sport-and-health:1}
Um coração maior, com volume sistólico maior e frequência de repouso
menor; mais capilares e \omterm{def:g10:cells-common-unit:organelle}{mitocôndrias} nas \omterm{def:g10:muscles-and-joints:muscle}{fibras musculares}; mais
hemácias; reservas maiores de \omterm{def:g10:exercise-and-energy:glycogen}{glicogênio} (e melhor regulação da
glicemia e da pressão).
\end{solution}

\begin{solution}{exo:g10:sport-and-health:2}
O volume sistólico aumentou, de modo que os mesmos \qty{5}{L/min} são
entregues em menos batimentos: $Q = f \times V_s$ com um $V_s$ maior
precisa de um $f$ menor.
\end{solution}

\begin{solution}{exo:g10:sport-and-health:3}
Cerca de uma hora por dia de atividade moderada. Não: caminhar, ir de
bicicleta à escola, escadas e tarefas domésticas contam --- o que
importa é o movimento que eleva o \omterm{def:g10:exercise-and-energy:expenditure}{gasto energético}.
\end{solution}

\begin{solution}{exo:g10:sport-and-health:4}
O uso de substâncias ou de métodos que elevam o desempenho atropelando
a regulação do corpo em vez de construir uma adaptação. O EPO atropela
o controle do número de hemácias; os estimulantes atropelam a fadiga, o
aviso que protege o coração.
\end{solution}

\begin{solution}{exo:g10:sport-and-health:5}
Fadiga persistente, frequência cardíaca de repouso em alta, desempenho
em queda (e ainda sono perturbado, infecções frequentes, dor em \omterm{def:g10:muscles-and-joints:muscle}{tendão}
ou em osso).
\end{solution}

\begin{solution}{exo:g10:sport-and-health:6}
Depois de 12 semanas: $72 - 60 = 12$ batimentos para quatro sessões,
$72 - 64 = 8$ para duas. A adaptação é proporcional à demanda --- mais
sessões, mais mudança --- e nos dois grupos ela desacelera à medida que
o novo estado é alcançado.
\end{solution}

\begin{solution}{exo:g10:sport-and-health:7}
$V_s = Q/f$: de $5000/72 \approx \qty{69}{mL}$ para $5000/54 \approx
\qty{93}{mL}$, um fator de $72/54 \approx 1.33$.
\end{solution}

\begin{solution}{exo:g10:sport-and-health:8}
$1 - 0.55 = 45\%$ menor. O primeiro degrau, de menos ativo para baixa
(1.00 para 0.80), traz o maior ganho.
\end{solution}

\begin{solution}{exo:g10:sport-and-health:9}
\qty{2.0}{kg}, ou seja $2.0/55 \approx 3.6\%$ da massa dela. O volume de
sangue dela cai, a frequência cardíaca dela sobe para manter o débito, a
temperatura do corpo dela sobe, o desempenho dela cai; as cãibras ficam
prováveis.
\end{solution}

\begin{solution}{exo:g10:sport-and-health:10}
$10 \times 1.1^7 \approx \qty{19.5}{km}$, de modo que os \qty{20}{km}
são alcançados na oitava semana: cerca de dois meses.
\end{solution}

\begin{solution}{exo:g10:sport-and-health:11}
Cada fibra engrossa acrescentando miofibrilas, e o sistema nervoso
aprende a recrutar mais fibras ao mesmo tempo e em melhor sincronia; as
duas coisas elevam a força. O número de fibras é fixo na vida adulta.
\end{solution}

\begin{solution}{exo:g10:sport-and-health:12}
O oxigênio por litro sobe por $56/45 \approx 1.24$, cerca de 24\%; com
o mesmo \omterm{def:g10:heart-lungs-effort:output}{débito cardíaco} e a mesma extração, o
$\dot V\!\mathrm{O_2}$max sobe até a mesma fração. Mas um sangue com
56\% de \omterm{def:g10:cells-common-unit:cell}{células} é muito mais espesso; desacelerado por um coração que
dorme, ele pode coagular num vaso coronário ou cerebral.
\end{solution}

\begin{solution}{exo:g10:sport-and-health:13}
Na altitude, o próprio corpo eleva o EPO dele em resposta a uma falta
real de oxigênio, numa quantidade controlada, e o reduz de novo quando
a falta acaba: a regulação está funcionando. O EPO injetado impõe um
número que a regulação nunca escolheria, sem falta que o justifique e
sem freio.
\end{solution}

\begin{solution}{exo:g10:sport-and-health:14}
O músculo em atividade capta glicose sem insulina, e o músculo treinado
permanece mais sensível a ela; caminhar reduz, portanto, a glicemia
diretamente e diminui a insulina necessária a cada refeição. Um
medicamento reduz o açúcar mas deixa a causa --- um músculo insensível
e sem uso --- no lugar, e não faz nada pelo coração, pelos ossos nem
pela massa corporal, que a caminhada também melhora.
\end{solution}

\begin{solution}{exo:g10:sport-and-health:15}
A adaptação é construída durante o descanso entre demandas repetidas e
gradualmente aumentadas: um mês de carga intensiva tem mais chance de
produzir excesso de treino e lesão do que boa forma, já que o \omterm{def:g10:muscles-and-joints:muscle}{tendão} e o
osso se adaptam ao longo de meses, e não de semanas. E os ganhos se
desmontam poucas semanas depois de parar: nada fica ``garantido'' por
um ano com um único arranque. A atividade moderada e regular ao longo
do ano faz o que o mês não consegue.
\end{solution}

\begin{solution}{pb:g10:sport-and-health:1}
\textbf{1.} $4800/72 \approx \qty{67}{mL}$ antes; $4800/56 \approx
\qty{86}{mL}$ depois.

\textbf{2.} $42 \times 62 \approx \qty{2.6}{L/min}$; $50 \times 62 =
\qty{3.1}{L/min}$: um ganho de 19\%.

\textbf{3.} Débito máximo $\dot V\!\mathrm{O_2}/0.150$: de \num{17.4}
para \qty{20.7}{L/min}; volume sistólico máximo a 200 batimentos: de
\num{87} para \qty{103}{mL}.

\textbf{4.} O coração maior e mais forte explica o volume sistólico;
mais capilares e \omterm{def:g10:cells-common-unit:organelle}{mitocôndrias} nas fibras explicariam um ganho de
extração.

\textbf{5.} Um fator de 4 em 26 semanas é uma média de $4^{1/26}
\approx 1.055$, cerca de 5.5\% por semana --- dentro da regra dos 10\%
($1.1^{26}
\approx 12$ teria permitido muito mais).

\textbf{6.} Excesso de treino: frequência cardíaca de repouso em alta,
sono perturbado, infecção, desempenho em queda, dor num tecido de
reparo lento.

\textbf{7.} Uma lesão por estresse do osso da canela (uma fratura por
estresse ou uma inflamação da superfície dele): o osso se adapta mais
devagar que todos os tecidos, e a carga foi dobrada em quinze dias.

\textbf{8.} A recuperação é incompleta de uma sessão para a seguinte; o
corpo está num estado de tensão persistente, e a adaptação que tinha
reduzido a frequência está sendo revertida.

\textbf{9.} Uma semana de quase repouso e de sono (a adaptação é
construída no descanso); depois duas semanas com metade da distância
anterior, sem trabalho de velocidade, reconstruindo um décimo por
semana; a canela avaliada por um médico e poupada até parar de doer
(escute a dor); beber e comer para o reparo.

\textbf{10.} Um estimulante: ele mascara a fadiga, o aviso que protege
o coração da exaustão e do superaquecimento.

\textbf{11.} Sam $32 \times 80 \approx \qty{2.6}{L/min}$; Alex $48
\times 64 \approx \qty{3.1}{L/min}$. O de Alex é maior mesmo em termos
absolutos; e a vida cotidiana --- escadas, caminhar, carregar o próprio
corpo --- move a massa do próprio corpo, de modo que o número por
quilograma é o que decide como isso é sentido.

\textbf{12.} Sam $12/32 \approx 38\%$ do máximo dele, Alex $12/48 =
25\%$. Sam, trabalhando acima de um terço do teto dele, chega sem
fôlego.

\textbf{13.} Um músculo sem uso capta pouca glicose do sangue e perde a
sensibilidade dele à insulina; as mesmas refeições exigem mais insulina
a cada ano e a glicemia começa a subir.

\textbf{14.} De 1.00 para 0.70: cerca de 30\% menos risco.

\textbf{15.} A \omterm{def:g10:sport-and-health:training}{atividade física} é qualquer movimento que eleve o \omterm{def:g10:exercise-and-energy:expenditure}{gasto
energético}: ir a pé ao trabalho, escadas em vez do elevador, tarefas
domésticas --- meia hora por dia disso não exige esporte nem
equipamento.

\textbf{16.} A 1\% ao ano, uma margem de 10\% são dez anos de queda.

\textbf{17.} O pico de Sam foi mais baixo e ele vem perdendo osso há
vinte e cinco anos; o punho dele cruzou o limiar em que uma queda o
quebra, enquanto o osso de Alex, construído sob carga na adolescência e
mantido desde então, tem dez anos de folga.

\textbf{18.} Uma adaptação é uma estrutura mantida pelo uso, e não uma
aquisição permanente: ela é desmontada quando a demanda para e
reconstruída quando ela volta.

\textbf{19.} Doenças do coração e das artérias, diabetes tipo 2,
osteoporose (e ainda alguns cânceres e a depressão). Diabetes: o
músculo inativo perde a sensibilidade dele à insulina, a glicemia sobe e
o pâncreas é sobrecarregado.

\textbf{20.} Frequência cardíaca de repouso 54 contra 78;
$\dot V\!\mathrm{O_2}$max \num{48} contra \qty{32}{mL/min} por
quilograma; cerca de três horas por semana de corrida, mantidas por
catorze anos.
\end{solution}
